\section*{Chapitre \ref{ch:pl:time-series-databases-and-the-alternatives} --- Bases de données de séries temporelles et leurs alternatives}

\begin{solution}{exo:pl:time-series-databases-and-the-alternatives:1}
Consultation ponctuelle~: $3.8 / 0.05$, environ 80 fois plus rapide. Coupe transversale~: $313 / 7$, plus de 40 fois plus lente. (Les facteurs exacts varient d'une
mesure à l'autre~; leur ordre de grandeur, lui, reste le même.)
\end{solution}

\begin{solution}{exo:pl:time-series-databases-and-the-alternatives:2}
L'index est un B-tree trié par (symbole, temps)~: la consultation descend quelques pages jusqu'à la dernière entrée du symbole
antérieure à l'instant demandé et lit une seule ligne. La coupe transversale a besoin de chaque ligne de chaque symbole~; l'index n'offre
aucun raccourci, et le stockage par lignes lit chaque ligne entière, une à une, pour n'en utiliser que deux champs.
\end{solution}

\begin{solution}{exo:pl:time-series-databases-and-the-alternatives:3}
$6.5 \times 60 = 390$ \omterm{def:pl:time-series-databases-and-the-alternatives:bucket}{intervalles de temps}. S007 n'a aucune cotation dans 31 d'entre eux~: un intervalle sans donnée ne produit aucune ligne, et une série de barres qui
a besoin de chaque minute doit combler explicitement les vides (en reportant le dernier mid, avec un compte nul).
\end{solution}

\begin{solution}{exo:pl:time-series-databases-and-the-alternatives:4}
La vitesse d'une machine dérive pendant un banc d'essai (fréquence, autres tâches, cache des pages). Lancer un moteur sept fois, puis le suivant,
attribuerait la dérive aux moteurs~; l'entrelacement expose chaque moteur aux mêmes conditions à chaque tour.
\end{solution}

\begin{solution}{exo:pl:time-series-databases-and-the-alternatives:5}
Les deux sont dominées par le coût fixe $c_0$ (ouverture du fichier du jour, lecture de son pied de fichier, planification). Et la requête ponctuelle, telle qu'écrite
(la dernière cotation à midi ou avant, triée par temps, limite à une), parcourt chaque cotation du symbole jusqu'à midi, soit à peu près la moitié de sa journée, et les trie,
alors que la plage ne touche que les 153 lignes de l'heure, élaguées par les \omterm{def:pl:columnar-formats:pushdown}{cartes de zones}. Ce sont les lignes touchées, et non les lignes renvoyées, qui fixent le
coût variable.
\end{solution}

\begin{solution}{exo:pl:time-series-databases-and-the-alternatives:6}
Coût égal lorsque $(1-f)\,3.8 + 313 f = (1-f)\,0.05 + 7 f$ (dans les unités des médianes mesurées)~: $f = (3.8-0.05)/[(3.8-0.05) +
(313-7)] = 1.2~\%$.
\end{solution}

\begin{solution}{exo:pl:time-series-databases-and-the-alternatives:7}
Les moteurs concordent sur les cinq formes~: une ligne pour la consultation ponctuelle, 171 cotations dans l'heure, 358 barres d'une minute, 2\,000 transactions jointes et
50 fourchettes moyennes.
\end{solution}

\begin{solution}{exo:pl:time-series-databases-and-the-alternatives:8}
Une seule forme de requête, celle du fournisseur, peut ne pas correspondre à la charge de travail de la firme~; un passage par moteur, en séquence, confond les moteurs avec le cache des pages
(le second passage trouve les données déjà en cache) et avec la dérive de la machine~; rien ne vérifiait que les deux donnaient la même réponse. Mesurer les formes de
requêtes propres à la firme sur ses propres données, avec des réponses vérifiées, des tours entrelacés, et des résultats à froid et à chaud présentés séparément.
\end{solution}

\begin{solution}{pb:pl:time-series-databases-and-the-alternatives:1}
\begin{enumerate}
\item Consultation ponctuelle (la dernière cotation avant un ordre), plage (l'heure d'un symbole), barres (barres d'une minute sur une journée), jointure as-of (chaque transaction avec sa cotation),
coupe transversale (fourchette moyenne par symbole sur un mois).
\item Consultation ponctuelle, plage et barres touchent peu de lignes~; la jointure as-of en touche deux tranches~; la coupe transversale touche toutes les lignes.
\item Données triées par symbole puis par temps, avec des \omterm{def:pl:columnar-formats:pushdown}{cartes de zones} (ou un index), de sorte que la fenêtre d'un symbole soit une tranche contiguë.
\item Les mêmes idées, conditionnées~: un serveur qui ingère, compacte, conserve et répond, avec des opérations propres au temps intégrées.
\item Remplacer les enregistrements d'un \omterm{def:pl:time-series-databases-and-the-alternatives:bucket}{intervalle de temps} par un résumé~; dans la \omterm{def:pl:capturing-and-storing-tick-data:tier}{politique de rétention}, les ticks sont purgés après quelques mois tandis que les barres sont conservées pendant des années.
\item Un moteur embarqué s'exécute dans le processus qui interroge (pas d'aller-retour, pas de sérialisation), mais renonce aux écritures concurrentes, au contrôle d'accès
et au service partagé~; un serveur offre tout cela, mais paie l'aller-retour.
\item Exécuter chaque opérateur sur un lot de valeurs d'une colonne, de sorte que le surcoût d'appel soit payé par lot et que les boucles internes soient des boucles
serrées sur des tableaux.
\item Faire circuler à travers les filtres les positions des lignes retenues, et ne lire les autres colonnes que pour les lignes qui survivent.
\item Un B-tree parcourt quelques pages jusqu'à une seule ligne~; un balayage de fichier paie encore l'ouverture, les métadonnées et le décodage d'un bloc.
\item Le plus rapide~: SQLite pour la consultation ponctuelle, numpy pour la plage, les barres et la coupe transversale, numpy et SQLite à égalité pour la jointure as-of~; le plus lent~:
pandas sur quatre formes, SQLite sur la coupe transversale.
\item 1, 153, 359, 2\,000 et 50~; car une réponse rapide et fausse est le résultat le plus courant d'un banc d'essai, et cinq moteurs qui ne s'accordent pas rendraient les temps
dénués de sens.
\item SQLite 0.05~ms et 313~ms, DuckDB 3.8~ms et 7~ms.
\item Il est écrit pour cette disposition et ces requêtes (tableaux triés, un décalage par symbole)~; chaque nouvelle question exige un nouveau code et de nouveaux tests, que
un moteur généraliste fournit gratuitement.
\item Les coûts de démarrage du moteur tout juste lancé et tout ce qui n'était pas encore en cache, que la première requête d'une session paie.
\item La machine, sa charge, l'état du cache des pages, la taille et la disposition des données, et les versions des moteurs.
\item \textbf{Résultat nommé.} Sur le mois de cotations, l'index de SQLite l'emporte sur la consultation ponctuelle (0.05~ms contre 3.8 pour DuckDB), les
moteurs en colonnes l'emportent sur les balayages (la coupe transversale en 7~ms contre 313 pour SQLite), un moteur écrit à la main bat les deux là où il s'applique,
et DuckDB devance SQLite sur une charge mixte dès que plus d'environ 1.2~\% des requêtes sont des coupes transversales.
\item DuckDB (ou Polars)~: la charge de travail porte sur des balayages et des agrégations sur l'historique.
\item Un index en processus (SQLite, ou un tableau trié en mémoire avec décalages, comme dans le moteur numpy), à côté du routeur~: une
consultation ponctuelle doit coûter des microsecondes, et non une ouverture de fichier et un plan.
\item Lorsque de nombreux écrivains et lecteurs partagent des données en direct, qu'il faut un contrôle d'accès et une rétention sans que la firme ait à les construire, et que
le banc d'essai sur les formes propres à la firme le montre gagnant.
\item La forme de la requête~: combien de lignes elle doit toucher, rapporté aux coûts fixes et par ligne du moteur.
\end{enumerate}
\end{solution}

\begin{solution}{iq:pl:time-series-databases-and-the-alternatives:1}
Une base de données optimisée pour des enregistrements indexés par le temps, ajoutés dans l'ordre temporel et lus par plages de temps, avec des \omterm{def:pl:time-series-databases-and-the-alternatives:bucket}{intervalles de temps}, des jointures as-of et
une rétention intégrée. Pour la recherche sur les \omterm{def:pl:capturing-and-storing-tick-data:tick}{données tick}, des fichiers dans des formats colonnaires ouverts et un moteur embarqué suffisent souvent~; un serveur est
utile lorsque de nombreux écrivains et lecteurs partagent des données en direct.
\iqlookfor{La catégorie relève du conditionnement d'idées connues~; c'est le mélange des requêtes qui tranche.}
\end{solution}

\begin{solution}{iq:pl:time-series-databases-and-the-alternatives:2}
Elle ne lit que les colonnes nécessaires, par lots, avec des boucles vectorisées serrées et la \omterm{def:pl:time-series-databases-and-the-alternatives:vectorised}{matérialisation tardive}, de sorte que les balayages coûtent peu par ligne.
Une ligne isolée coûte toujours un montant fixe (ouverture, plan, décodage d'un bloc), qu'un index sur un stockage par lignes évite.
\iqlookfor{Coût fixe face au coût par ligne.}
\end{solution}

\begin{solution}{iq:pl:time-series-databases-and-the-alternatives:3}
Utiliser ses propres données et ses formes de requêtes~; vérifier d'abord que les deux renvoient des réponses identiques~; entrelacer les passages~; présenter à part les résultats à froid et à chaud
avec les dispersions~; préciser la machine, la charge et le cache~; faire varier la taille des données pour voir comment chacun évolue.
\iqlookfor{Justesse avant vitesse~; entrelacement~; froid et chaud.}
\end{solution}

\begin{solution}{iq:pl:time-series-databases-and-the-alternatives:4}
Lorsque les données sont des fichiers qu'une seule équipe lit (une \omterm{def:pl:a-tick-store-on-open-formats:store}{base de ticks}, des jeux de données de recherche) et que les requêtes sont analytiques~: pas de serveur à faire tourner, pas
d'aller-retour, copie sans sans copie vers les dataframes. Un serveur est nécessaire pour de nombreuses écritures concurrentes, pour un contrôle d'accès ou pour un seul jeu de données
partagé et en direct.
\iqlookfor{Qui écrit, et combien de lecteurs le partagent.}
\end{solution}

\begin{solution}{iq:pl:time-series-databases-and-the-alternatives:5}
Inventorier les requêtes auxquelles il répond (formes, fréquences, besoins de latence) et ses écrivains~; construire des candidats sur des formats ouverts (fichiers
et moteurs embarqués, un ou deux serveurs)~; les soumettre à un banc d'essai sur les données de la firme, avec des réponses vérifiées par rapport à l'ancien système~;
migrer les lecteurs forme par forme, en gardant les deux systèmes jusqu'à ce que les réponses concordent pendant des semaines.
\iqlookfor{Inventaire de la charge de travail, banc d'essai vérifié, exploitation en parallèle.}
\end{solution}
